\section{Fekete convergence when subadditivity begins on a tail}

\begin{lemma}\label{lem:tail-subadditive-arithmetic}
\leanok
\lean{EnvelopeProperty.TailSubadditive.apply_mul_add_le}
Let $u:\mathbf N\to\mathbf R$ and $N\in\mathbf N$ satisfy
$u_{m+n}\leq u_m+u_n$ for $m,n\geq N$.
If $n,r\geq N$, then $u_{kn+r}\leq ku_n+u_r$ for every $k\geq0$.
\end{lemma}
\begin{proof}
The case $k=0$ is equality. In the induction step apply tail subadditivity
to $n$ and $kn+r$, which are both at least $N$, and then use the induction
hypothesis. No value before the specified tail is used.
\end{proof}

\begin{lemma}\label{lem:tail-subadditive-eventual-upper-bound}
\leanok
\lean{EnvelopeProperty.TailSubadditive.eventually_div_lt_of_div_lt}
Under the same assumptions, let $n\geq N$, $n>0$. If $u_n/n<L$, then
$u_p/p<L$ for every sufficiently large $p$.
\uses{lem:tail-subadditive-arithmetic}
\end{lemma}
\begin{proof}
There are finitely many residue classes modulo $n$. For each $0\leq r<n$
put $s=Nn+r\geq N$. The previous lemma gives
$u_{kn+s}/(kn+s)\leq(ku_n+u_s)/(kn+s)$ once the denominator is positive.
The right side tends to $u_n/n<L$. Translating $k$ by $N$ gives the desired
eventual estimate on the whole residue class $np+r$. Taking a common bound
for the finitely many classes proves the assertion.
\end{proof}

\begin{theorem}\label{thm:tail-subadditive-fekete}
\leanok
\lean{EnvelopeProperty.TailSubadditive.tendsto_limit}
Suppose additionally that $u_n\geq0$ for every $n\geq N$.
\lean{EnvelopeProperty.TailSubadditive.limit}
Then $u_n/n$ converges to the finite nonnegative real number
$\inf\{u_n/n:n\geq\max(N,1)\}$.
\uses{lem:tail-subadditive-eventual-upper-bound}
\end{theorem}
\begin{proof}
The displayed set is nonempty and bounded below by zero, so its real
infimum exists and is nonnegative. For every sufficiently large index,
the normalized value lies in this set and hence is above its infimum.
Given a real $L$ strictly above the infimum, some member is smaller than
$L$. Its positive index lies beyond $N$, so the preceding lemma bounds
all sufficiently large normalized values above by $L$. These two estimates
are precisely convergence to the infimum in the order topology.
\end{proof}
